% Tapis de Quillen, with Quillen's offprint and preprint (cote 162-5, typed part): archive cote 162-5, site page 47 (PDF page 48).
% typescript. Draft compared: none.
% AI-transcribed by Claude Opus 5.5 (claude-opus-5-5, adaptive thinking, effort high), Message Batches API, batch t162_full (msgbatch_01UAh3zyWPd454AeDkhnmQ33), 2026-11-13; not yet checked by eye.
\scanpage{162-5/47}
\deleted{tés, les morphismes \uline{toutes} les correspondances \inserted[typed]{co}homologiques (en vert\unsure{u}}
\deleted{du thorème de Thom rappelé plus haut). On trouve donc un foncteur}
\deleted{\unsure{pleinemeht|pleinement} fidèle de cette c atégorie dans celle d es vectoriels}
\deleted{gradués \xout{\unsure{dexd}} positivement de dim. finie sur $\underline{F}_2$ , qui est déjà abélien}
\deleted{ne, et si on veut rendre comme plus haut $\underline{F}_2(-1)$ inversible, on trouve}
\deleted{tous les vectoriels gradués de dim. finie sur $\underline{F}_2$ . C'est un peu trop}
\deleted{trivial comme structure, et on peut faire mieux, \xout{en prenant comme}}
\inserted[typed]{Pour trouver une catégorie universelle,}
\xout{\unsure{correspondances. Pour ceci}} on associe à toute X \xout{\unsure{le sous-anneau ensemble}}
\xout{\unsure{de}} l'ensemble \inserted[typed]{S(X)} des sous-variétés \xout{S(X)} de X, modulo équivalence homolo%
gique. Donc S(X) \hand{$\subset$} $H^{\cdot}(X)$, et on prouve qu'on trouve en fait un sous-%
anneau de $H^{\cdot}(X)$ (grâce à un théorème de Thom qui caractérise S(X) en
termes d'opérations de Steenrod, cf plus bas). On prouve de même
(par le même théorème) que les S(XxY) sont stables par composition de\unsure{s}
correspondances cohomologiques( cela ne semble pas trivial, puisque
les S(X) ne sont pas stables par Gysin...). D'oû une catégorie addi-%
tive dont les objets sont ceux de C, les flêches sont les S(XxY).
Il semble qu'elle doive jouer le rôle universel que j'ai envisagé
plus haut. En tous cas, sa définition géom\xout{m}étrique est fort simple !
Nous allons en donner une interpétation algébrique. Pour ce\unsure{ci},
considérons l'anneau des opérations cohomologiques \inserted[typed]{``stables''} \deleted{\xout{additives sur}} \deleted{i.e.}
\margin{vertical line in the left margin, from this line down to the end of the typed text}\deleted{\inserted[typed]{des endomorphismes du}}
\margin{\unsure{$\xi$}}\deleted{\xout{le} foncteur X \quad $H^{\cdot}$(X) qui sont additifs, et qui en degré zéro se}
\deleted{réduisent à un multiple de l'identité.} C'est un anneau gradué, avec
application diagonale anticommutative, \inserted[typed]{i.e. commutative (car. 2 !)} associative et unitaire. Son
dual gradué est une algèbre \inserted[typed]{graduée anti\unsure{c}commutative donc \xout{commutative},} commutative,avec application diagonale
associative et unitaire; il est réduit à $\underline{F}_2$ en degré zéro. Donc il
corespond à un schéma en groupes \unsure{aff}ine G sur $\underline{F}_2$ , qui est d'ailleurs
le foncteur des \unsure{$\otimes$}-automorphismes \inserted[typed]{additifs} du foncteur $H^{\cdot}$ sur C \inserted{qui \unsure{opèrent}
trivialement sur \deleted{\unsure{l'objet}} $H^1(S^1)$. \deleted{\illegible{1}}}

\mspage{47}
